\chapter{抓取的接触运动学与力/形封闭理论}
\label{ch:grasp-closure}

% ════════════════════════════════════════════════════════════════
%  机械臂建模、规划与控制部 · 抓取理论章
%  地面真值：Murray–Li–Sastry《MLS》Ch5（多指手运动学）+ Lynch–Park《MR》§12.1–12.2。
%  按编写规范：经典正典“基于原书内容重述”——教材化、深入、记号归一，非翻译非抄录。
%  §一.8 自包含铁律：核心定义/定理/判据/推导完整写入正文，\cite 只标出处。
%  记号归一（preamble 宏 \SO \SE \so \se；本书 wrench/twist 力在前 [\rho;\phi]）：
%    - se(3) 坐标“平移在前” \xi=[\rho;\phi]（见 notation.tex）；wrench 同序 w=[f;m]（力在前）。
%    - MLS/MR 原文 wrench 多取“力矩在前” F=(m,f)；本章统一转“力在前”，转换处显式标注。
%    - 抓取映射 G 与 cooperative_manipulation/dual_multi_arm 一致：w=G f，G_i=[I;\skewm{r}_i]。
%    - 摩擦锥/GIWC/Minkowski–Weyl 直接 \cref contact_mechanics，不重推。
%  姊妹章“抓取规划”（力/形封闭判据 → 抓点采样/打分/学习）尚未建章，本章 §延伸阅读以散文承接、
%    不 \cref 不存在锚点（grep 已确认 ch:grasp-planning 当前不存在）。
% ════════════════════════════════════════════════════════════════

% —— 本章局部数学算子（章文件 \input 入正文，故用 \providecommand+\operatorname 兜底，
%    与 cooperative_manipulation.tex 一致，避免与 preamble 既有定义冲突）——
\providecommand{\rank}{\operatorname{rank}}
\providecommand{\range}{\operatorname{range}}
\providecommand{\rowsp}{\operatorname{row}}
\providecommand{\spanop}{\operatorname{span}}
\providecommand{\pos}{\operatorname{pos}}
\providecommand{\conv}{\operatorname{conv}}
\providecommand{\intop}{\operatorname{int}}
% 反对称（叉积）矩阵：接触点位置向量 r 的 \widehat{r}（与 \cref{ch:lie} 的 (·)^\wedge 同义）
\providecommand{\skewm}[1]{\widehat{#1}}

本章把前面建立的\textbf{抓取映射} $\mathbf{G}$（\cref{sec:cm-grasp} 的 $\mathbf{w}=\mathbf{G}\mathbf{f}$）与\textbf{接触摩擦锥}（\cref{def:friction-cone}）合流，回答一个工程上最朴素也最深刻的问题：\emph{一组手指（或机器人末端）按住一个物体，到底“按住了没有”？} 答案分两个层次——纯几何上能否封住物体的全部运动自由度（\textbf{形封闭}，form closure），以及借助摩擦后能否抵抗任意外力旋量（\textbf{力封闭}，force closure）。围绕这两个判据，我们先补齐\emph{接触一阶运动学}（法/切向、roll/slide/break 三模式、多接触枚举），再给出\emph{接触力旋量基}与\emph{手部雅可比}这两件连接“手”与“物”的工具，最后把封闭性化为线性代数与线性规划上的可判定问题。

% ════════════════════════════════════════════════════════════════
\begin{practice}[前置自测]
开始之前，先试答下列问题。能流畅作答可速读本章表格与速查卡；若卡住（答不出 $\ge 2$ 题），正是本章要为你解决的。
\begin{enumerate}
  \item 一个夹爪两指压住一个立方体，你能说清“稳”这个字的\emph{数学}含义吗？它和“能抵抗重力”、“能抵抗任意方向的推搡”是不是一回事？
  \item 无摩擦时，一根手指对物体只能施加哪个方向的力？平面上要“锁死”一个物体（任何方向都推不动）最少需要几根这样的手指？空间里呢？
  \item 接触点处“滚动”“滑动”“分离”三种情形，分别对应接触点相对速度的什么约束？哪一种会让接触失效？
  \item 把摩擦考虑进来后，为什么空间里\emph{三根}带摩擦的手指就可能锁死物体，而无摩擦却要七根？多出来的“摩擦力”到底买到了什么？
  \item 给定抓取映射 $\mathbf{G}$ 和各接触摩擦锥 $\mathrm{FC}$，“能抵抗任意外力”如何写成一个可计算的判据？它和“原点落在某个凸包内部”是同一件事吗？
  \item 手指自身关节够不够用，会不会影响封闭性？“物体能被这只手任意搬动”和“这只手能锁死物体”是不是对偶的两面？
\end{enumerate}
\noindent\emph{答案线索}：1、5 是本章\cref{sec:gc-forceclosure} 力封闭的核心；2 是\cref{sec:gc-formclosure} 形封闭的接触数定理；3 回\cref{sec:gc-contactkin} 一阶接触运动学；4 是 Nguyen 三接触定理（\cref{thm:gc-nguyen}）；6 回\cref{sec:gc-handjac} 手部雅可比的可抓取性 $\range(\mathbf{G}^\top)\subseteq\range(\mathbf{J}_h)$。本章每步推导都建立在\cref{ch:lie}（旋量/伴随）、\cref{ch:contact-mechanics}（摩擦锥/GIWC）与\cref{ch:mr-coopforce}（抓取映射）之上。
\end{practice}

\section*{本章目标}
学完本章，你应当能够：
\begin{enumerate}
  \item \textbf{建立}单接触一阶运动学：用接触法向 $\hat{\mathbf{n}}$ 写不可穿透约束 $\mathcal{F}^\top(\mathcal{V}_A-\mathcal{V}_B)\ge0$，并据相对速度把接触判为滚动 R / 滑动 S / 分离 B 三模式；
  \item \textbf{枚举}多接触下物体的可行运动旋量集（多面体凸锥），理解“$j$ 个接触约束正张成 wrench 空间 $\Leftrightarrow$ 物体被完全约束”；
  \item \textbf{掌握}三类接触模型的力旋量基 $\mathbf{B}_{c_i}$——无摩擦点接触（$1$ 维）、有摩擦点接触 PCWF（$3$ 维）、软指 SF（$4$ 维）——及其摩擦锥 $\mathrm{FC}$ 是闭凸锥；
  \item \textbf{推导}手部雅可比 $\mathbf{J}_h$ 与基本抓取约束 $\mathbf{J}_h\dot{\boldsymbol{\theta}}=\mathbf{G}^\top\mathcal{V}$，并用 $\range(\mathbf{G}^\top)\subseteq\range(\mathbf{J}_h)$ 判可抓取（manipulable）、用 $\ker\mathbf{J}_h^\top$ 识别结构力；
  \item \textbf{严格定义}力封闭 $\mathbf{G}(\mathrm{FC})=\mathbb{R}^p$，并掌握其三条等价判据（接触 wrench 正张成全空间 / 原点在 wrench 凸包内部 / 无分离超平面），及与摩擦锥、GIWC 的关系；
  \item \textbf{严格定义}形封闭（纯几何无摩擦）、陈述一阶形封闭的平面 $\ge4$ / 空间 $\ge7$ 接触数定理，并用线性规划判定；
  \item \textbf{对照}形封闭与力封闭的异同，知道何时用哪一个、它们各自的代价与盲区。
\end{enumerate}

\subsection*{本章知识导航}
本章是一条“\emph{从局部接触约束逐级累积成全局封闭判据}”的主线。结构如下：
\begin{enumerate}
  \item \cref{sec:gc-contactkin}——\textbf{接触一阶运动学}：单接触的不可穿透/激活约束、roll-slide-break 三模式、多接触可行旋量锥。这是“运动被约束了多少”的几何骨架。
  \item \cref{sec:gc-wrenchbasis}——\textbf{接触力模型与力旋量基}：无摩擦/PCWF/软指三模型的 $\mathbf{B}_{c_i}$ 与摩擦锥 $\mathrm{FC}$。这是“力能施加多少”的对偶骨架。
  \item \cref{sec:gc-handjac}——\textbf{手部雅可比与基本抓取约束}：把手指关节速度接进来，$\mathbf{J}_h\dot{\boldsymbol{\theta}}=\mathbf{G}^\top\mathcal{V}$，可抓取性与结构力。
  \item \cref{sec:gc-forceclosure}——\textbf{力封闭}：$\mathbf{G}(\mathrm{FC})=\mathbb{R}^p$、三等价判据、Nguyen 三接触定理、与 GIWC 的统一。
  \item \cref{sec:gc-formclosure}——\textbf{形封闭}：纯几何定义、接触数定理、LP 判据、二阶效应的告诫。
  \item \cref{sec:gc-compare}——\textbf{形封闭 vs 力封闭对比}：一表收束全章。
\end{enumerate}

% ════════════════════════════════════════════════════════════════
\section{接触一阶运动学：滚动、滑动、分离}
\label{sec:gc-contactkin}

\paragraph{动机：一只手能不能锁住物体，先看运动而非看力。}
研究封闭性有两条入手路径。一条从\emph{力}入手（手指能施加哪些力、合起来能否对抗外扰），是\cref{sec:gc-forceclosure} 的力封闭；另一条从\emph{运动}入手——把摩擦先搁置，纯看“几何上接触把物体的哪些瞬时运动堵死了”，剩下没堵死的运动若为零，物体就被锁死。后者是本节的\textbf{一阶接触运动学}，也是形封闭的地基。它有一个迷人的好处：\emph{只用接触点与接触法向，不需要任何力学或摩擦参数}，是最“便宜”的分析。

\paragraph{反面：忽略接触模式会得出物理上不可能的运动。}
若把接触简单当成“两点重合”的等式约束，会犯两类错误。其一，\textbf{把单边当双边}：接触是单边的——手指能推不能拉（除非粘住），物体可以\emph{离开}手指（分离），却不能\emph{压进}手指（穿透）。用等式约束 $\mathbf{p}_A=\mathbf{p}_B$ 会同时禁止分离，于是算出“物体被锁死”的假结论。其二，\textbf{混淆滚动与滑动}：滚动接触切向无相对速度（决定能否用切向摩擦），滑动接触切向有相对滑移（摩擦力方向被滑移方向钉死）。把二者混为一谈，后续摩擦力建模与稳定性判断全错。本节用一阶（$\dot d=0$）分析把这三种模式精确区分开。

\subsection{单接触：距离函数与不可穿透约束}
\label{sec:gc-contact-single}

设两刚体 $A,B$ 的位形分别为局部坐标列向量 $\mathbf{q}_1,\mathbf{q}_2$，合并写 $\mathbf{q}=(\mathbf{q}_1,\mathbf{q}_2)$。定义\textbf{距离函数} $d(\mathbf{q})$：两体分离时 $d>0$、相切时 $d=0$、互相穿透时 $d<0$\cite{lynch2017modern}。当 $d>0$ 时两体各有六自由度、互不约束；当 $d=0$ 时，是否保持接触由 $d$ 的各阶时间导数决定。对一条轨迹 $\mathbf{q}(t)$，前两阶导数为
\begin{equation}
  \dot d=\frac{\partial d}{\partial\mathbf{q}}\dot{\mathbf{q}},
  \qquad
  \ddot d=\dot{\mathbf{q}}^\top\frac{\partial^2 d}{\partial\mathbf{q}^2}\dot{\mathbf{q}}+\frac{\partial d}{\partial\mathbf{q}}\ddot{\mathbf{q}}.
  \label{eq:gc-d-derivs}
\end{equation}
梯度 $\partial d/\partial\mathbf{q}$ 编码接触法向信息；二阶项 $\partial^2 d/\partial\mathbf{q}^2$ 编码相对曲率。

\begin{definition}[一阶接触分析]
\label{def:gc-first-order}
\textbf{一阶接触分析}假定接触处只有法向信息 $\partial d/\partial\mathbf{q}$ 可用，曲率及更高阶几何未知，分析截断在 \cref{eq:gc-d-derivs} 的一阶项：当 $\dot d=0$ 时即认为两体保持接触。其余各阶（$\ddot d$ 等）的可行/不可行判定见下表（按字典序逐阶判断，遇第一个非零导数定胜负）。
\end{definition}

\begin{table}[H]\centering\small
\caption{接触是否维持：按距离函数各阶导数逐级判定（首个非零项决定）}
\label{tab:gc-contact-derivatives}
\begin{tabular}{cccccl}
\toprule
$d$ & $\dot d$ & $\ddot d$ & $\cdots$ & & 结论 \\
\midrule
$>0$ & & & & & 无接触 \\
$<0$ & & & & & 不可行（穿透） \\
$=0$ & $>0$ & & & & 接触中，但正在分离 \\
$=0$ & $<0$ & & & & 不可行（穿透） \\
$=0$ & $=0$ & $>0$ & & & 接触中，但正在分离 \\
$=0$ & $=0$ & $<0$ & & & 不可行（穿透） \\
$=0$ & $=0$ & $=0$ & $\cdots$ & & （需更高阶） \\
\bottomrule
\end{tabular}
\end{table}

\noindent 只有当\emph{所有}各阶导数为零，接触才被精确维持。一阶分析只看 $d=\dot d=0$，故可能把实际“正在分离/穿透”的二阶情形误判为“维持”——这一保守性正是后文形封闭判据“一阶充分、非必要”的根源（\cref{sec:gc-formclosure}）。

\subsection{用接触法向改写：不可穿透与激活约束}
\label{sec:gc-impenetrability}

显式距离函数往往难得，但接触法向总是可知的。设 $\hat{\mathbf{n}}\in\mathbb{R}^3$ 为接触法向单位向量（世界系，约定指向物体 $A$ 内部），$\mathbf{p}_A,\mathbf{p}_B\in\mathbb{R}^3$ 为接触点在 $A,B$ 上的世界系表示。两点初始重合，但其速度 $\dot{\mathbf{p}}_A,\dot{\mathbf{p}}_B$ 可不同。则 $\dot d=0$ 的接触维持条件写为
\begin{equation}
  \hat{\mathbf{n}}^\top(\dot{\mathbf{p}}_A-\dot{\mathbf{p}}_B)=0,
  \label{eq:gc-dot-pn}
\end{equation}
而不可穿透（$\dot d\ge0$，两体可分离不可压入）为 $\hat{\mathbf{n}}^\top(\dot{\mathbf{p}}_A-\dot{\mathbf{p}}_B)\ge0$。

把它写成运动旋量（twist）形式。设 $A,B$ 的空间运动旋量为 $\mathcal{V}_A=(\mathbf{v}_A,\boldsymbol{\omega}_A)$、$\mathcal{V}_B=(\mathbf{v}_B,\boldsymbol{\omega}_B)$（本书 twist 平移在前，\cref{def:so3-se3}），刚体上接触点速度为 $\dot{\mathbf{p}}_A=\mathbf{v}_A+\boldsymbol{\omega}_A\times\mathbf{p}_A=\mathbf{v}_A+\skewm{\boldsymbol{\omega}}_A\mathbf{p}_A$，对 $B$ 同理。再引入沿接触法向施加单位力对应的\textbf{接触法向力旋量}（contact-normal wrench）。

\begin{note}[排序对照：本书 wrench 力在前，与 MR/MLS 力矩在前互为置换]
《MR》§12.1 与《MLS》Ch5 把 wrench 写成“力矩在前” $\mathcal{F}=(\mathbf{m},\mathbf{f})$，于是单位法向力旋量记 $\mathcal{F}=(\mathbf{p}_A\times\hat{\mathbf{n}},\,\hat{\mathbf{n}})=(\skewm{\mathbf{p}}_A\hat{\mathbf{n}},\,\hat{\mathbf{n}})$。本书 wrench 取\emph{力在前} $\mathcal{F}=[\mathbf{f};\mathbf{m}]$（与运动旋量 $[\boldsymbol{\rho};\boldsymbol{\phi}]$ 同序，\nameref{ch:notation}），故同一物理量在本书序里写成
\begin{equation}
  \mathcal{F}=\begin{bmatrix}\hat{\mathbf{n}}\\ \skewm{\mathbf{p}}_A\hat{\mathbf{n}}\end{bmatrix}\in\mathbb{R}^6
  \quad(\text{力在前；下半块 }\skewm{\mathbf{p}}_A\hat{\mathbf{n}}=\mathbf{p}_A\times\hat{\mathbf{n}}\ \text{是力矩}).
  \label{eq:gc-normal-wrench}
\end{equation}
与之配套，twist 也取力学功配对 $\mathcal{F}^\top\mathcal{V}=\mathbf{f}^\top\mathbf{v}+\mathbf{m}^\top\boldsymbol{\omega}$。引文公式凡“力矩在前”者，本章一律转置到本书序后再用，转换处显式标注，结论（功、对偶、张成性）与排序无关。
\end{note}

借力旋量与运动旋量的功配对，不可穿透与激活约束写成不依赖坐标选取的紧凑形式：
\begin{align}
  \text{（不可穿透约束）}\quad & \mathcal{F}^\top(\mathcal{V}_A-\mathcal{V}_B)\ge0,
  \label{eq:gc-impenetrability}\\
  \text{（激活约束，一阶接触维持）}\quad & \mathcal{F}^\top(\mathcal{V}_A-\mathcal{V}_B)=0.
  \label{eq:gc-active}
\end{align}
当 $B$ 为静止固定件（$\mathcal{V}_B=\mathbf{0}$），简化为 $\mathcal{F}^\top\mathcal{V}_A\ge0$：若 $\mathcal{F}^\top\mathcal{V}_A>0$ 称 $\mathcal{F}$ 与 $\mathcal{V}_A$ \textbf{互斥}（repelling，接触分离）；若 $\mathcal{F}^\top\mathcal{V}_A=0$ 称\textbf{互易}（reciprocal，约束激活、接触维持）。

\subsection{三模式：滚动 R、滑动 S、分离 B}
\label{sec:gc-rsb}

满足激活约束 \cref{eq:gc-active} 的运动旋量称\textbf{一阶 roll-slide 运动}——接触维持，但还分滚动与滑动两子类。

\begin{definition}[接触模式 R/S/B]
\label{def:gc-rsb}
对维持中的单点接触：
\begin{itemize}
  \item \textbf{滚动（R）}：两体在接触点处无相对运动，即接触点速度相等
  \begin{equation}
    \dot{\mathbf{p}}_A=\mathbf{v}_A+\skewm{\boldsymbol{\omega}}_A\mathbf{p}_A=\mathbf{v}_B+\skewm{\boldsymbol{\omega}}_B\mathbf{p}_B=\dot{\mathbf{p}}_B.
    \label{eq:gc-rolling}
  \end{equation}
  （含两体相对静止、“粘连/sticking”的情形。）
  \item \textbf{滑动（S）}：满足激活约束 \cref{eq:gc-active} 但\emph{不}满足滚动约束 \cref{eq:gc-rolling}——接触维持却有切向滑移。
  \item \textbf{分离（B，breaking free）}：满足不可穿透 \cref{eq:gc-impenetrability} 但\emph{不}满足激活约束（$\mathcal{F}^\top(\mathcal{V}_A-\mathcal{V}_B)>0$）——接触正在断开。
\end{itemize}
$n$ 个接触的系统模式记为各接触标签的串接；三标签下 $n$ 接触最多 $3^n$ 种模式，部分因运动学不相容而不可行。
\end{definition}

\begin{note}[为什么滚动是“无滑移”而不是“有转动”]
“滚动”在日常语义里强调“滚”（有转动），但在接触运动学里它\emph{只}指接触点相对速度为零（\cref{eq:gc-rolling}），与有无整体转动无关——两体完全相对静止（“粘住”）也是滚动的特例。区分 R/S 的唯一标准是接触界面切向\emph{有没有相对滑移}，因为这决定了切向摩擦力能否“咬住”：滚动时摩擦力可取摩擦锥内任意方向（静摩擦），滑动时摩擦力被滑移方向反向钉死（动摩擦，\cref{def:friction-cone} 的锥面上）。这一区分在\cref{ch:contact-mechanics} 的接触力学里反复出现。
\end{note}

\begin{example}[平面单接触的三模式划分]
\label{ex:gc-single-contact}
设 $B$ 静止、$A$ 限于 $z=0$ 平面运动，接触点 $\mathbf{p}_A=[1,2,0]^\top$、法向 $\hat{\mathbf{n}}=[0,1,0]^\top$\cite{lynch2017modern}。本书序的法向力旋量 \cref{eq:gc-normal-wrench} 的非零分量为 $f_y=1$、$m_z=(\mathbf{p}_A\times\hat{\mathbf{n}})_z=1$。取 $A$ 的平面运动旋量 $\mathcal{V}_A=[v_{Ax},v_{Ay},\omega_{Az}]$，激活约束 \cref{eq:gc-active} 给出
\[
  v_{Ay}+\omega_{Az}=0,
\]
这是 $(\omega_{Az},v_{Ax},v_{Ay})$ 空间里过原点的一个平面：平面之上（$v_{Ay}+\omega_{Az}>0$）是分离 B，之下是穿透（不可行）。再叠加滚动约束 \cref{eq:gc-rolling}（本例化简为 $v_{Ax}-2\omega_{Az}=0$ 且 $v_{Ay}+\omega_{Az}=0$）得平面内一条直线：直线上是滚动 R，平面上其余 roll-slide 点是滑动 S。一张图里 B/R/S 三区一目了然——这正是后文“形封闭 = 把可行旋量锥压成原点”的局部素材。
\end{example}

\subsection{多接触：可行运动旋量的多面体凸锥}
\label{sec:gc-multicontact}

现设物体 $A$ 受 $n$ 个接触（来自 $m$ 个其他物体，$n\ge m$）。每个接触 $i$ 把 $\mathcal{V}_A$ 约束到其六维旋量空间的一个半空间，边界是五维超平面 $\mathcal{F}_i^\top\mathcal{V}_A=\mathcal{F}_i^\top\mathcal{V}_{j(i)}$。取所有约束之交，得\textbf{可行运动旋量多面体}
\begin{equation}
  \mathcal{V}=\{\mathcal{V}_A\mid \mathcal{F}_i^\top(\mathcal{V}_A-\mathcal{V}_{j(i)})\ge0,\ i=1,\dots,n\},
  \label{eq:gc-twist-polytope}
\end{equation}
其中 $\mathcal{F}_i$ 为第 $i$ 个接触法向力旋量（指向物体 $A$）。若某接触的半空间约束不改变 $\mathcal{V}$ 则称该约束\textbf{冗余}。一般 $\mathcal{V}$ 可有六维内部（无约束激活）、五维面（一约束激活）、……直至零维点；$\mathcal{V}_A$ 落在 $k$ 维面上意味着有 $6-k$ 个独立（非冗余）接触约束激活。

\begin{insight}[全部接触体静止 $\Rightarrow$ 可行旋量集是“扎根原点的锥”]
当所有施约束的物体都静止（$\mathcal{V}_j=\mathbf{0}$），每个约束超平面 \cref{eq:gc-impenetrability} 都过原点（称\emph{齐次约束}），可行集退化为扎根原点的\textbf{齐次多面体凸锥}
\begin{equation}
  \mathcal{V}=\{\mathcal{V}_A\mid \mathcal{F}_i^\top\mathcal{V}_A\ge0,\ i=1,\dots,n\}.
  \label{eq:gc-twist-cone}
\end{equation}
此时一个判据呼之欲出：\emph{若 $\{\mathcal{F}_i\}$ 正张成（positively span）六维 wrench 空间——等价地其凸包内部含原点——则 $\mathcal{V}$ 缩成单点 $\{\mathbf{0}\}$，静止接触完全锁死物体，即形封闭}（\cref{sec:gc-formclosure}）。注意这里出现了“一组 wrench 正张成全空间”这个核心条件，它将在力封闭里以\emph{完全相同}的形式再次登场——形/力封闭的数学内核是同一个正张成问题，差别只在“用哪一组 wrench”（纯法向 vs 摩擦锥棱）。
\end{insight}

% ════════════════════════════════════════════════════════════════
\section{接触力模型与力旋量基 \texorpdfstring{$\mathbf{B}_{c_i}$}{B}}
\label{sec:gc-wrenchbasis}

\paragraph{承上：从“运动被约束”到“力能施加”。}
\cref{sec:gc-contactkin} 回答了接触把运动堵死了多少；本节回答其对偶面——接触能向物体\emph{施加哪些力旋量}。二者由同一接触法向 $\hat{\mathbf{n}}$ 联系：\emph{能约束运动的方向，恰是能施力的方向}（\cref{eq:gc-impenetrability} 的 $\mathcal{F}$ 同时是约束法向与可施力向）。一个接触能施加的力旋量被两件东西刻画：一个\textbf{力旋量基} $\mathbf{B}_{c_i}$（指明哪些维度可施力）与一个\textbf{摩擦锥} $\mathrm{FC}_{c_i}$（指明这些维度上的可行幅值集合）。

\subsection{接触坐标系与三类接触模型}
\label{sec:gc-contact-models}

约定第 $i$ 个接触坐标系 $C_i$ 的 $z$ 轴沿接触处\emph{指向物体内的曲面法向}\cite{murray1994mathematical}。手指在 $C_i$ 处施加的接触力（在 $C_i$ 系下表达）记 $\mathbf{f}_{c_i}$。常见三模型：

\begin{definition}[无摩擦点接触 / 有摩擦点接触 PCWF / 软指 SF]
\label{def:gc-contact-models}
\hfill
\begin{enumerate}
  \item \textbf{无摩擦点接触}（frictionless point contact，$1$ 维）：只能沿法向施正压力，
  \begin{equation}
    \mathbf{B}_{c_i}=\begin{bmatrix}0\\0\\1\\0\\0\\0\end{bmatrix}\in\mathbb{R}^{6\times1},
    \qquad \mathrm{FC}_{c_i}=\{f_z\in\mathbb{R}: f_z\ge0\}.
    \label{eq:gc-Bc-frictionless}
  \end{equation}
  \item \textbf{有摩擦点接触}（point contact with friction, \textbf{PCWF}，$3$ 维）：可施法向力与切向摩擦力，受库仑摩擦锥约束，
  \begin{equation}
    \mathbf{B}_{c_i}=\begin{bmatrix}\mathbf{I}_3\\ \mathbf{0}_3\end{bmatrix}\in\mathbb{R}^{6\times3},
    \quad
    \mathrm{FC}_{c_i}=\Big\{\mathbf{f}\in\mathbb{R}^3:\ \sqrt{f_x^2+f_y^2}\le\mu f_z,\ f_z\ge0\Big\}.
    \label{eq:gc-Bc-pcwf}
  \end{equation}
  \item \textbf{软指接触}（soft-finger，\textbf{SF}，$4$ 维）：PCWF 之上再加一个绕法向的扭转力矩 $m_z$（受扭转摩擦系数 $\gamma$ 限），
  \begin{equation}
    \mathbf{B}_{c_i}=\begin{bmatrix}\mathbf{I}_3 & \mathbf{0}\\ \mathbf{0} & \mathbf{e}_3\end{bmatrix}\in\mathbb{R}^{6\times4},
    \ \,
    \mathrm{FC}_{c_i}=\Big\{(\mathbf{f},m_z):\sqrt{f_x^2+f_y^2}\le\mu f_z,\ |m_z|\le\gamma f_z,\ f_z\ge0\Big\},
    \label{eq:gc-Bc-sf}
  \end{equation}
  其中 $\mathbf{e}_3=[0,0,1]^\top$ 选出绕法向的力矩分量。
\end{enumerate}
$\mathbf{B}_{c_i}\in\mathbb{R}^{p\times m_i}$ 的列数 $m_i$ 是该接触可独立施加的力的维度；空间抓取 $p=6$，平面抓取 $p=3$。
\end{definition}

\noindent 摩擦锥的几何（半顶角 $\alpha=\arctan\mu$、二阶锥的凸性）已在\cref{def:friction-cone}、\cref{prop:cone-convex} 详述，此处直接复用（\cref{eq:friction-cone-contact}）；线性化为多面锥（pyramidal）以入 LP/QP 的做法见\cref{sec:contact-linearize}。三模型的“单点力维 $\to$ 抓取矩阵尺寸 $\to$ 内力维度”系统分类已在\cref{tab:cm-contact-models} 给出，本章不重列。

\begin{definition}[摩擦锥须满足的凸锥公理]
\label{def:gc-fc-axioms}
任一接触的摩擦锥 $\mathrm{FC}_{c_i}\subset\mathbb{R}^{m_i}$ 须满足：（1）$\mathrm{FC}_{c_i}$ 是 $\mathbb{R}^{m_i}$ 的\emph{闭}子集且\emph{内部非空}；（2）锥性与凸性——$\mathbf{f}_1,\mathbf{f}_2\in\mathrm{FC}_{c_i}\Rightarrow\alpha\mathbf{f}_1+\beta\mathbf{f}_2\in\mathrm{FC}_{c_i}$ 对一切 $\alpha,\beta\ge0$。该接触\emph{实际}能施加的力旋量集合为 $\{\mathbf{B}_{c_i}\mathbf{f}_{c_i}:\mathbf{f}_{c_i}\in\mathrm{FC}_{c_i}\}$。
\end{definition}

\begin{note}[“能施力的方向”与“能约束运动的方向”是同一组]
\cref{sec:gc-impenetrability} 里无摩擦点接触的运动约束写作 $[0,0,1,0,0,0]\,\mathcal{V}_{f,c}=0$（$z$ 向相对速度为零）——其行向量恰是无摩擦点接触力旋量基 $\mathbf{B}_{c_i}$ 的转置。一般地，接触约束 $\mathbf{B}_{c_i}^\top\mathcal{V}_{f,c}=\mathbf{0}$ 把“手指相对物体的速度”限制在“可施力方向的正交补”里：\emph{凡能施力之处，运动即被锁；凡不能施力之处（如无摩擦切向、点接触绕法向转），运动自由}。这条“施力方向 $=$ 约束方向”的对偶，是下一节手部雅可比推导的出发点。
\end{note}

\subsection{把接触力搬到物体系：回到抓取映射 \texorpdfstring{$\mathbf{G}$}{G}}
\label{sec:gc-graspmap-recall}

单个接触的力旋量经伴随变换搬到物体质心系，叠加 $N$ 个接触即得物体合力旋量 $\mathbf{w}_o=\mathbf{G}\mathbf{f}_c$，其中第 $i$ 块 $\mathbf{G}_i=\mathrm{Ad}_{g_{oc_i}^{-1}}^\top\mathbf{B}_{c_i}$（把 $C_i$ 系力旋量映成物体系力旋量）。对 PCWF 点接触（每点 $3$ 维力），这正是\cref{eq:cm-Gi} 的
\begin{equation}
  \mathbf{G}_i=\begin{bmatrix}\mathbf{I}_3\\ \skewm{\mathbf{r}}_i\end{bmatrix}\in\mathbb{R}^{6\times3},
  \qquad \mathbf{w}_o=\mathbf{G}\mathbf{f}_c=\sum_{i=1}^N\mathbf{G}_i\mathbf{f}_{c_i},
  \label{eq:gc-graspmap}
\end{equation}
$\mathbf{r}_i$ 为接触点相对质心的位置（本书序力在前，上块 $\mathbf{I}_3$ 搬力、下块 $\skewm{\mathbf{r}}_i$ 经力臂生力矩）。$\mathbf{G}$ 的完整推导、零空间 $\dim\ker\mathbf{G}=cN-6$、内力/外力分解 $\mathbf{f}=\mathbf{G}^+\mathbf{w}_{\text{ext}}+(\mathbf{I}-\mathbf{G}^+\mathbf{G})\boldsymbol{\eta}$ 均见\cref{ch:mr-coopforce}（\cref{eq:cm-grasp}、\cref{eq:cm-decomp}），多臂/双臂对偶见\cref{ch:mr-multiarm}（\cref{eq:ma-grasp}），本章直接取用。

\begin{definition}[抓取的代数描述]
\label{def:gc-grasp}
一个\textbf{抓取}由抓取映射 $\mathbf{G}\in\mathbb{R}^{p\times m}$（$m=\sum_i m_i$）与复合摩擦锥 $\mathrm{FC}=\mathrm{FC}_{c_1}\times\cdots\times\mathrm{FC}_{c_N}\subset\mathbb{R}^m$ 完全刻画；$\mathrm{FC}$ 满足\cref{def:gc-fc-axioms} 的闭凸锥公理。物体可被施加的力旋量集合为
\begin{equation}
  \{\mathbf{w}_o=\mathbf{G}\mathbf{f}_c:\ \mathbf{f}_c\in\mathrm{FC}\}=\mathbf{G}(\mathrm{FC}).
  \label{eq:gc-G-FC}
\end{equation}
\end{definition}

% ════════════════════════════════════════════════════════════════
\section{手部雅可比 \texorpdfstring{$\mathbf{J}_h$}{Jh} 与基本抓取约束}
\label{sec:gc-handjac}

\paragraph{承上：把手指自身的关节接进来。}
至此 $\mathbf{G}$ 描述了“接触力 $\to$ 物体力旋量”，却把手指当成了能在接触点任意施力的“理想接触”。真实手指是\emph{有限自由度的关节链}，其末端只能产生雅可比像空间内的速度与力。本节引入\textbf{手部雅可比} $\mathbf{J}_h$，把关节速度接到接触约束上，得到连接“手关节世界”与“物体世界”的\textbf{基本抓取约束}，并由此引出两个对偶性质——\emph{可抓取性}（手能否实现物体任意运动）与\emph{结构力}（手关节无法控制、由结构吸收的力）。

\subsection{接触约束与手部雅可比}
\label{sec:gc-Jh-derive}

设第 $i$ 根手指有 $n_i$ 个关节、关节角 $\boldsymbol{\theta}_i\in\mathbb{R}^{n_i}$，指尖相对接触系的物体运动旋量为 $\mathcal{V}_{f_i,c_i}$。接触约束 \cref{eq:gc-active} 的一般形式（对力旋量基 $\mathbf{B}_{c_i}$ 的所有可施力方向）为\cite{murray1994mathematical}
\begin{equation}
  \mathbf{B}_{c_i}^\top\,\mathcal{V}_{f_i,c_i}=\mathbf{0}
  \qquad(\text{指尖相对物体的速度，须正交于一切可施力方向}).
  \label{eq:gc-contact-constraint}
\end{equation}
用指尖运动学把 $\mathcal{V}_{f_i,c_i}$ 拆成“手指关节贡献”减“物体运动贡献”，并用伴随变换归算到 $C_i$ 系，可化为（推导见\cref{der:gc-Jh}）
\begin{equation}
  \mathbf{B}_{c_i}^\top\,\mathrm{Ad}_{g_{c_if_i}}^{-1}\mathbf{J}^{s}_{f_i}(\boldsymbol{\theta}_i)\,\dot{\boldsymbol{\theta}}_i
  =\big(\mathrm{Ad}_{g_{oc_i}^{-1}}^\top\mathbf{B}_{c_i}\big)^\top\mathcal{V}
  =\mathbf{G}_i^\top\mathcal{V},
  \label{eq:gc-contact-row}
\end{equation}
其中 $\mathbf{J}^{s}_{f_i}$ 是第 $i$ 指指尖的（空间）雅可比，$\mathcal{V}$ 是物体运动旋量。

\begin{derivation}[基本抓取约束的逐行化简]
\label{der:gc-Jh}
约束 \cref{eq:gc-contact-constraint} 的左端是“指尖相对物体在 $C_i$ 系的速度”。把它写成三段串联：指尖相对手指基座（$=\mathbf{J}^{s}_{f_i}\dot{\boldsymbol{\theta}}_i$，手指正运动学的雅可比，\cref{def:ak-jacobian}）、手指基座相对手掌（常量）、手掌相对物体（$=\mathcal{V}$，物体运动）。三段用伴随变换 $\mathrm{Ad}$ 归算到公共的 $C_i$ 系（\cref{sec:adjoint}）并相减，左乘 $\mathbf{B}_{c_i}^\top$ 投到“可施力方向”，再用伴随的群同态性质 $\mathrm{Ad}_{g_1g_2}=\mathrm{Ad}_{g_1}\mathrm{Ad}_{g_2}$ 把物体侧整理成 $\mathbf{G}_i^\top\mathcal{V}$，即得 \cref{eq:gc-contact-row}。关键一步是\emph{物体侧的 $\mathbf{B}_{c_i}^\top\mathrm{Ad}_{g_{oc_i}^{-1}}$ 恰是 $\mathbf{G}_i^\top$}——这把手的约束与抓取映射 $\mathbf{G}$ 在同一坐标语言下接上了。\qed
\end{derivation}

把 $N$ 个接触的 \cref{eq:gc-contact-row} 竖直堆叠，定义\textbf{手部雅可比}
\begin{equation}
  \mathbf{J}_h(\boldsymbol{\theta},\mathbf{X}_o)
  =\begin{bmatrix}
    \mathbf{B}_{c_1}^\top\mathrm{Ad}_{g_{c_1f_1}}^{-1}\mathbf{J}^{s}_{f_1} & & \\
    & \ddots & \\
    & & \mathbf{B}_{c_N}^\top\mathrm{Ad}_{g_{c_Nf_N}}^{-1}\mathbf{J}^{s}_{f_N}
  \end{bmatrix}\in\mathbb{R}^{m\times n},
  \label{eq:gc-Jh-def}
\end{equation}
$n=\sum_i n_i$ 为总关节数，$m=\sum_i m_i$。于是诸约束合写为\textbf{基本抓取约束}：

\begin{theorem}[基本抓取约束]
\label{thm:gc-grasp-constraint}
多指手抓持刚体时，手指关节速度 $\dot{\boldsymbol{\theta}}\in\mathbb{R}^n$ 与物体运动旋量 $\mathcal{V}\in\mathbb{R}^p$ 满足
\begin{equation}
  \boxed{\ \mathbf{J}_h(\boldsymbol{\theta},\mathbf{X}_o)\,\dot{\boldsymbol{\theta}}=\mathbf{G}^\top\mathcal{V}.\ }
  \label{eq:gc-grasp-constraint}
\end{equation}
其物理含义：等式两端都是“接触点处、相对接触系、沿可施力方向的速度”——左端由手指关节经 $\mathbf{J}_h$ 产生，右端由物体运动经 $\mathbf{G}^\top$ 分发到各接触点（\cref{eq:cm-vel-dual} 的多指版）。一个抓取由 $(\mathbf{J}_h,\mathbf{G},\mathrm{FC})$ 三件套完整刻画。
\end{theorem}

\begin{note}[\texorpdfstring{$\mathbf{J}_h$}{Jh} 与 \texorpdfstring{$\mathbf{G}^\top$}{GT} 的角色不可互换]
易混点：\cref{eq:gc-grasp-constraint} 左右都带“接触点速度”，但来源不同。$\mathbf{G}^\top\mathcal{V}$ 是\emph{物体若以 $\mathcal{V}$ 运动、接触点应有的速度}（运动学对偶，纯几何）；$\mathbf{J}_h\dot{\boldsymbol{\theta}}$ 是\emph{手指关节若以 $\dot{\boldsymbol{\theta}}$ 运动、指尖在接触点产生的速度}（手指运动学）。基本抓取约束要求二者在“可施力方向”上相等——否则指尖与物体在接触点撕开或压入。把 $\mathbf{G}^\top$ 误当 $\mathbf{J}_h$（或反之）会得到维度都对、物理全错的方程：$\mathbf{G}$ 只含接触几何（$\mathbf{r}_i$），$\mathbf{J}_h$ 含手指构型（$\boldsymbol{\theta}$）。
\end{note}

\subsection{可抓取性与结构力：两个对偶判据}
\label{sec:gc-manip-structural}

基本抓取约束 \cref{eq:gc-grasp-constraint} 的代数结构直接给出两个互为对偶的能力判据。

\begin{definition}[可抓取/可操作抓取 manipulable grasp]
\label{def:gc-manipulable}
若对物体的\emph{任意}运动旋量 $\mathcal{V}$，都存在手指关节速度 $\dot{\boldsymbol{\theta}}$ 满足 \cref{eq:gc-grasp-constraint}，则称该抓取在当前构型 $(\boldsymbol{\theta},\mathbf{X}_o)$ 下\textbf{可抓取}（可操作）。
\end{definition}

\begin{proposition}[可抓取性的秩判据]
\label{prop:gc-manipulable}
抓取在 $(\boldsymbol{\theta},\mathbf{X}_o)$ 可抓取，当且仅当
\begin{equation}
  \range(\mathbf{G}^\top)\subseteq\range(\mathbf{J}_h).
  \label{eq:gc-manipulable}
\end{equation}
即“物体运动经 $\mathbf{G}^\top$ 落入的接触速度”都能被手指雅可比像空间实现。
\end{proposition}
\begin{proof}
可抓取 $\Leftrightarrow$ 对任意 $\mathcal{V}$，$\mathbf{G}^\top\mathcal{V}\in\range(\mathbf{J}_h)$（这样 \cref{eq:gc-grasp-constraint} 才有解 $\dot{\boldsymbol{\theta}}$）。$\mathcal{V}$ 遍历全空间时 $\mathbf{G}^\top\mathcal{V}$ 遍历 $\range(\mathbf{G}^\top)$，故条件即 $\range(\mathbf{G}^\top)\subseteq\range(\mathbf{J}_h)$。$\square$
\end{proof}

\noindent 可抓取性\emph{不}要求 $\mathbf{J}_h$ 单射：当 $\mathbf{J}_h$ 行满秩且列多于行（$n>m$）时有非平凡零空间，凡 $\dot{\boldsymbol{\theta}}_N\in\ker\mathbf{J}_h$ 都使接触点不动——称\textbf{内部运动}（internal motion），可叠加到手指速度上而不改变物体运动，使“指-物系统”成为冗余机械臂。值得注意：即便每根手指本身不冗余，$\mathbf{B}_{c_i}^\top$ 的“屏蔽”作用（滤掉不违反接触约束的关节速度，如绕点接触法向的自转）也会带来这种冗余。

对偶地，从力侧看（接触力与关节力矩由 $\boldsymbol{\tau}=\mathbf{J}_h^\top\mathbf{f}_c$ 联系，\cref{def:ak-jacobian} 的静力学对偶）：

\begin{proposition}[结构力 structural forces]
\label{prop:gc-structural}
落在 $\ker\mathbf{J}_h^\top$ 中的接触力 $\mathbf{f}_c$ 不需任何关节力矩即可被手指结构平衡，称\textbf{结构力}：它们由连杆/关节的机械结构（而非电机）吸收，手电机对其\emph{无法控制}。故\emph{可由手指主动控制}的接触力受限于 $\range(\mathbf{J}_h)$ 之内；力封闭分析中“假定一切 $\mathrm{FC}$ 内的力都可施加”这一前提，须 $\mathbf{J}_h$ 行满秩（无不可控结构力维度）方成立。
\end{proposition}

\begin{insight}[可抓取性与力封闭：手够不够 vs 接触够不够]
两个能力判据各管一段，别混。\textbf{力封闭}（\cref{sec:gc-forceclosure}）问“\emph{接触几何 $\mathbf{G}$ 与摩擦锥 $\mathrm{FC}$}够不够抵抗任意外力”——只看 $\mathbf{G},\mathrm{FC}$，把手指当理想施力器。\textbf{可抓取性}问“\emph{手指关节 $\mathbf{J}_h$}够不够实现物体任意运动 / 施加任意接触力”——只看 $\mathbf{J}_h$。一个抓取可以接触几何完美（力封闭）却手指自由度不足（不可抓取，如三指压住物体但每指只一关节、无法独立调三维接触力），也可以反之。完整“好抓取”需\emph{二者兼备}加上结构力可控（$\mathbf{J}_h$ 行满秩），这是把理想 $\mathbf{G}$ 分析落到真实手的桥。
\end{insight}

\subsection{贯通范例：两 SCARA 手指抓盒}
\label{sec:gc-scara}

\paragraph{把 $\mathbf{G}$、$\mathbf{J}_h$、可抓取性、内部运动一次走通。}
前几节的 $\mathbf{G}$、$\mathbf{B}_{c_i}$、$\mathbf{J}_h$ 与可抓取性判据，单独看都抽象。这里用一个完整数值范例把它们\emph{串起来}\cite{murray1994mathematical}：两根 SCARA 手指各以\emph{软指接触}（SF，\cref{def:gc-contact-models} 第 3 类）夹住一个盒子的两侧，接触点在物体系下为 $\mathbf{p}_{c_1}=(0,-r,0)$、$\mathbf{p}_{c_2}=(0,+r,0)$（\cref{fig:gc-scara}）。我们将依次算出抓取映射 $\mathbf{G}$、分块手部雅可比 $\mathbf{J}_h$，并判定：该抓取\textbf{力封闭}（接触几何 + 软指锥够），却\textbf{不可抓取}（手指自由度不够实现物体某些运动），并解出其\textbf{内部运动}。这正是 \cref{ins:gc-scara}（亦即上节 \cref{def:gc-manipulable} 后强调的）“力封闭 $\ne$ 可抓取”最干净的实证。

\begin{figure}[ht]
\centering
\begin{tikzpicture}[scale=1.0,>=Stealth,font=\scriptsize]
% box
\draw[thick, main!60!black, fill=main!6] (-0.7,-0.5) rectangle (0.7,0.5);
\node at (0,0) {盒子};
\draw[<->] (-0.7,0.72)--(0.7,0.72); \node[above] at (0,0.7) {$2r$};
% contact points and normals
\fill (-0.7,0) circle (1.4pt) node[left=1pt]{$C_1$};
\fill (0.7,0) circle (1.4pt) node[right=1pt]{$C_2$};
\draw[->,second!70!black,thick] (-0.7,0)--(-0.25,0); % normal into box
\draw[->,second!70!black,thick] (0.7,0)--(0.25,0);
% left SCARA finger (vertical post + two links reaching in)
\draw[line width=2pt, gray!55!black] (-2.6,-0.9)--(-2.6,0.6); % post S1
\node[below] at (-2.6,-0.9) {$S_1$};
\draw[line width=1.8pt, third!60!black] (-2.6,0.3)--(-1.7,0.15)--(-0.7,0);
\fill[third!70!black] (-1.7,0.15) circle (1.3pt);
\node at (-2.15,0.45){$l_1$}; \node at (-1.15,0.32){$l_2$};
% right SCARA finger
\draw[line width=2pt, gray!55!black] (2.6,-0.9)--(2.6,0.6); \node[below] at (2.6,-0.9) {$S_2$};
\draw[line width=1.8pt, third!60!black] (2.6,0.3)--(1.7,0.15)--(0.7,0);
\fill[third!70!black] (1.7,0.15) circle (1.3pt);
\node at (2.15,0.45){$l_3$}; \node at (1.15,0.32){$l_4$};
% base distance
\draw[<->] (-2.6,-1.15)--(2.6,-1.15); \node[below] at (0,-1.15) {$2b$};
\end{tikzpicture}
\caption{两 SCARA 手指抓盒。手指立柱 $S_1,S_2$ 相距 $2b$，各经两段水平连杆（长 $l_1,l_2$ 与 $l_3,l_4$）伸到盒子两侧的软指接触 $C_1,C_2$（盒宽 $2r$）。每根手指是 4 自由度 SCARA：三个绕竖直 $z$ 轴的转动关节（前两个定位、第三个腕转）$+$ 一个沿 $z$ 的升降移动关节。}
\label{fig:gc-scara}
\end{figure}

\paragraph{第一步：抓取映射 $\mathbf{G}$（由接触系经伴随搬到物体系）。}
两接触系（$z$ 轴沿指向物体内的法向，\cref{sec:gc-contact-models}）相对物体系的姿态与位置为\cite{murray1994mathematical}
\begin{equation}
\mathbf{R}_{c_1}=\begin{bmatrix}0&1&0\\0&0&1\\1&0&0\end{bmatrix},\ \mathbf{p}_{c_1}=\begin{bmatrix}0\\-r\\0\end{bmatrix};\quad
\mathbf{R}_{c_2}=\begin{bmatrix}1&0&0\\0&0&-1\\0&1&0\end{bmatrix},\ \mathbf{p}_{c_2}=\begin{bmatrix}0\\r\\0\end{bmatrix}.
\label{eq:gc-scara-frames}
\end{equation}
软指力旋量基 $\mathbf{B}_{c_i}\in\mathbb{R}^{6\times4}$ 取 \cref{eq:gc-Bc-sf}（前三列法/切向力、第四列绕法向扭矩）。把它经伴随转置搬到物体系（\cref{eq:gc-graspmap} 的 $\mathbf{G}_i=\mathrm{Ad}_{g_{oc_i}^{-1}}^\top\mathbf{B}_{c_i}$，本书力在前序的 wrench 变换为 $\big[\begin{smallmatrix}\mathbf{R}_{c_i}&\mathbf{0}\\ \skewm{\mathbf{p}}_{c_i}\mathbf{R}_{c_i}&\mathbf{R}_{c_i}\end{smallmatrix}\big]$）。以手指 1 为例逐块算：上 $3$ 行 $1$--$3$ 列 $=\mathbf{R}_{c_1}$、第 $4$ 列 $=\mathbf{0}$；下 $3$ 行 $1$--$3$ 列 $=\skewm{\mathbf{p}}_{c_1}\mathbf{R}_{c_1}=\big[\begin{smallmatrix}-r&0&0\\0&0&0\\0&r&0\end{smallmatrix}\big]$、第 $4$ 列 $=\mathbf{R}_{c_1}\mathbf{e}_3=(0,1,0)^\top$，故
\begin{equation}
\mathbf{G}_1=\mathrm{Ad}_{g_{oc_1}^{-1}}^\top\mathbf{B}_{c_1}=
\begin{bmatrix}0&1&0&0\\0&0&1&0\\1&0&0&0\\-r&0&0&0\\0&0&0&1\\0&r&0&0\end{bmatrix}.
\label{eq:gc-scara-G1}
\end{equation}
对手指 2 同法，拼出抓取映射 $\mathbf{G}=[\mathbf{G}_1\ \mathbf{G}_2]\in\mathbb{R}^{6\times8}$：
\begin{equation}
\mathbf{G}=\begin{bmatrix}
0&1&0&0&\ 1&0&0&0\\
0&0&1&0&\ 0&0&-1&0\\
1&0&0&0&\ 0&1&0&0\\
-r&0&0&0&\ 0&r&0&0\\
0&0&0&1&\ 0&0&0&-1\\
0&r&0&0&\ -r&0&0&0
\end{bmatrix}.
\label{eq:gc-scara-G}
\end{equation}
逐行消元可验 $\rank\mathbf{G}=6$（六行线性无关），故 $\mathbf{G}$ 行满秩、\emph{满射}——配合软指锥含严格内力（“对夹挤压”），该抓取\textbf{力封闭}（满射 $+$ 严格内力，判据 \cref{prop:gc-internal-necessity}，下节详述）。\emph{力封闭只取决于接触几何 $\mathbf{G}$ 与摩擦锥，与具体用什么手指无关}——这一点下面与可抓取性形成鲜明对照。

\paragraph{第二步：SCARA 手指雅可比与分块手部雅可比 $\mathbf{J}_h$。}
取指尖长 $l_3=0$（接触即在 SCARA 末端），单根 SCARA 手指的空间雅可比为（本书 twist 平移在前，列依次为三转动 $+$ 一移动关节）\cite{murray1994mathematical}
\begin{equation}
\mathbf{J}^s_{f}=\begin{bmatrix}
0 & l_1c_1 & l_1c_1+l_2c_{12} & 0\\
0 & l_1s_1 & l_1s_1+l_2s_{12} & 0\\
0 & 0 & 0 & 1\\
0 & 0 & 0 & 0\\
0 & 0 & 0 & 0\\
1 & 1 & 1 & 0
\end{bmatrix},
\label{eq:gc-scara-Js}
\end{equation}
其中 $c_1=\cos\theta_1$、$c_{12}=\cos(\theta_1+\theta_2)$ 等。把它经 $\mathbf{B}_{c_i}^\top\mathrm{Ad}_{g_{c_if_i}}^{-1}(\cdot)$ 归算到接触系（\cref{eq:gc-Jh-def}），并\emph{特取 \cref{fig:gc-scara} 的对称构型}（$\mathbf{R}_{po}=\mathbf{I}$、$\mathbf{p}_{po}=(0,0,a)$），得分块手部雅可比 $\mathbf{J}_h=\mathrm{diag}(\mathbf{J}_{11},\mathbf{J}_{22})\in\mathbb{R}^{8\times8}$（手指 1 用 $l_1,l_2,\theta_1,\theta_2$；手指 2 用 $l_3,l_4,\theta_3,\theta_4$；$b$ 为半基距、$r$ 半盒宽）：
\begin{equation}
\mathbf{J}_{11}=\begin{bmatrix}
0&0&0&1\\
-b+r & -b+r+l_1c_1 & -b+r+l_1c_1+l_2c_{12} & 0\\
0 & l_1s_1 & l_1s_1+l_2s_{12} & 0\\
0&0&0&0
\end{bmatrix},\quad
\mathbf{J}_{22}=\begin{bmatrix}
b-r & b-r+l_3c_3 & b-r+l_3c_3+l_4c_{34} & 0\\
0&0&0&1\\
0 & -l_3s_3 & -l_3s_3-l_4s_{34} & 0\\
0&0&0&0
\end{bmatrix}.
\label{eq:gc-scara-Jh}
\end{equation}
\textbf{关键观察}：$\mathbf{J}_{11},\mathbf{J}_{22}$ 的\emph{第 4 行恒为零}。该行对应软指接触的\emph{绕法向扭转}速度分量（$\mathbf{B}_{c_i}^\top$ 第 4 行选出）——SCARA 三个转动关节都绕\emph{竖直} $z$ 轴，无法在\emph{水平}接触法向上产生扭转，故 $\mathbf{J}_h$ 第 $4,8$ 行整体为零。

\paragraph{第三步：判可抓取性——\emph{不}可抓取。}
取物体绕 $y$ 轴的转动旋量 $\mathcal{V}=(\mathbf{0};\,0,1,0)$（绕过两接触的 $y$ 轴），经 $\mathbf{G}^\top$ 分发到接触点（\cref{eq:gc-grasp-constraint} 右端）：
\begin{equation}
\mathbf{G}^\top\mathcal{V}=(\text{$\mathbf{G}$ 第 5 行})^\top=(0,0,0,1,\ 0,0,0,-1)^\top.
\label{eq:gc-scara-notmanip}
\end{equation}
此向量要求手指 1 的第 $4$ 分量（扭转速度）为 $1$、手指 2 的第 $4$（即整向量第 $8$）分量为 $-1$。但 $\mathbf{J}_h$ 第 $4,8$ 行为零，故 $\mathbf{G}^\top\mathcal{V}\notin\range(\mathbf{J}_h)$——\cref{eq:gc-grasp-constraint} 无解 $\dot{\boldsymbol{\theta}}$。由 \cref{prop:gc-manipulable}（$\range(\mathbf{G}^\top)\subseteq\range(\mathbf{J}_h)$ 不成立），该抓取在此构型\textbf{不可抓取}：手指\emph{够不着}“把盒子绕 $y$ 轴转一点”这一物体运动，因为这需要在接触处提供绕法向的扭转滑移，而竖直轴 SCARA 指做不到。

\paragraph{第四步：解内部运动。}
$\mathbf{J}_h$ 分块对角，每块 $\mathbf{J}_{ii}\in\mathbb{R}^{4\times4}$ 第 4 行为零、其余三行一般独立，故 $\rank\mathbf{J}_{ii}=3$、$\dim\ker\mathbf{J}_{ii}=1$，合得 $\dim\ker\mathbf{J}_h=2$。在 \cref{fig:gc-scara} 的对称构型下，手指 1 的指尖\emph{正对}接触点（水平伸直），即 $l_1s_1+l_2s_{12}=0$（指尖在指系内的 $y$ 向到达为零）且 $l_1c_1+l_2c_{12}=b-r$（$x$ 向到达恰为基到接触距），代入 \cref{eq:gc-scara-Jh} 知 $\mathbf{J}_{11}$ 的\emph{第 3 列恰为零}；对手指 2 同理第 3 列为零。于是内部运动为
\begin{equation}
\dot{\boldsymbol{\theta}}_{N_1}=(0,0,1,0,\ 0,0,0,0)^\top,\qquad
\dot{\boldsymbol{\theta}}_{N_2}=(0,0,0,0,\ 0,0,1,0)^\top,
\label{eq:gc-scara-internal}
\end{equation}
即\emph{各自转动该手指的第三关节（腕转）}——它使指尖\emph{绕接触点自转}而\emph{不移动物体}（$\mathbf{J}_h\dot{\boldsymbol{\theta}}_{N_i}=\mathbf{0}$，接触点速度为零）。这恰是 \cref{sec:gc-manip-structural} 所述“$\mathbf{J}_h$ 列多于秩 $\Rightarrow$ 指-物系统冗余”的具体兑现：两根原本各 4 自由度、不冗余的手指，因软指接触\emph{屏蔽}掉绕法向扭转（$\mathbf{B}_{c_i}^\top$ 滤波），反而获得 2 维内部运动。（更一般构型 $l_3\neq0$ 下内部运动仍存，但三个转动关节都有非零瞬时速度。）

\begin{insight}[一个抓取，力封闭却不可抓取：两判据正交]\label{ins:gc-scara}
本例把 \cref{def:gc-manipulable} 后强调的“力封闭 $\ne$ 可抓取”\emph{摆成了数字}：同一抓取，\textbf{力封闭}（$\mathbf{G}$ 满射 $+$ 软指内力，第一步）与\textbf{不可抓取}（$\range(\mathbf{G}^\top)\not\subseteq\range(\mathbf{J}_h)$，第三步）\emph{同时}成立，毫不矛盾——因为二者问的是\emph{不同}的事：力封闭只看“接触几何 $\mathbf{G}$ $+$ 摩擦锥能否抗任意外力”（与手指构造无关），可抓取性只看“手指雅可比 $\mathbf{J}_h$ 能否实现物体任意运动 / 调任意接触力”。竖直轴 SCARA 指撑得住盒子（力封闭），却\emph{转不动}盒子绕 $y$ 轴（不可抓取），因为它在接触处缺一个绕水平法向的扭转自由度。\textbf{记住这条边界}：抓取规划既要 $\mathbf{G},\mathrm{FC}$ 力封闭，又要 $\mathbf{J}_h$ 行满秩可抓取（且结构力可控，\cref{prop:gc-structural}），缺一不可。本例的两维内部运动则提示：富余的手指自由度可调内力、避奇异，是冗余资源而非负担。
\end{insight}

% ════════════════════════════════════════════════════════════════
\section{力封闭 Force Closure}
\label{sec:gc-forceclosure}

\paragraph{动机：抓得稳的数学定义。}
“抓稳了”最有用的数学定义是：\emph{无论外界对物体施加什么力旋量，手指都能在摩擦锥内调出接触力把它顶回去}。能做到这点的抓取称\textbf{力封闭}。它是抓取规划打分、灵巧操作可行性、装配能否夹持的统一底座，也是\cref{ch:contact-mechanics} 站立平衡 GIWC 判据在“抓取”语境下的孪生兄弟。

\subsection{严格定义与第一判据}
\label{sec:gc-fc-def}

\begin{definition}[力封闭抓取]
\label{def:gc-force-closure}
对一个抓取 $(\mathbf{G},\mathrm{FC})$，若对\emph{任意}外加力旋量 $\mathbf{w}_e\in\mathbb{R}^p$，都存在接触力 $\mathbf{f}_c\in\mathrm{FC}$ 使
\begin{equation}
  \mathbf{G}\mathbf{f}_c=-\mathbf{w}_e,
  \label{eq:gc-fc-def}
\end{equation}
则称该抓取为\textbf{力封闭}（force closure）。即接触力总能合成出与任意外扰相反的物体力旋量。
\end{definition}

\begin{proposition}[力封闭的像集判据]
\label{prop:gc-fc-image}
抓取力封闭，当且仅当
\begin{equation}
  \mathbf{G}(\mathrm{FC})=\mathbb{R}^p.
  \label{eq:gc-fc-image}
\end{equation}
即接触力经 $\mathbf{G}$ 能张满整个 $p$ 维力旋量空间（空间 $p=6$，平面 $p=3$）。
\end{proposition}
\begin{proof}
直接由\cref{def:gc-force-closure}：对任意 $\mathbf{w}_e$ 存在 $\mathbf{f}_c\in\mathrm{FC}$ 使 $\mathbf{G}\mathbf{f}_c=-\mathbf{w}_e$，等价于 $-\mathbf{w}_e$ 遍历 $\mathbb{R}^p$ 时总落在 $\mathbf{G}(\mathrm{FC})$ 内，即 $\mathbf{G}(\mathrm{FC})\supseteq\mathbb{R}^p$；又 $\mathbf{G}(\mathrm{FC})\subseteq\mathbb{R}^p$ 显然，故 $\mathbf{G}(\mathrm{FC})=\mathbb{R}^p$。$\square$
\end{proof}

\begin{note}[“力封闭”不等于“物体不动”]
力封闭只说\emph{接触摩擦锥能生成任意 wrench}，\emph{不}保证物体在某外力下静止不动。是否真不动还取决于\textbf{内力}够不够大：例如两指夹三角形、$\mu=1$、两指“互相看得见”（视线在双方摩擦锥内）是力封闭，但若电机出力不足、或只能朝固定方向出力，三角形仍可能在重力下滑落\cite{lynch2017modern}。力封闭是“能力”（capability），“顶得住”还需配套的力分配与内力调节（\cref{eq:cm-decomp} 的内力 $\mathbf{f}_{\text{int}}$）。这正呼应\cref{ch:mr-coopforce} 反复强调的“封闭性 $\ne$ 已实现的稳定平衡”。
\end{note}

\subsection{内力的必要性}
\label{sec:gc-fc-internal}

\begin{definition}[内力与严格内力]
\label{def:gc-internal}
若 $\mathbf{f}_N\in\ker\mathbf{G}\cap\mathrm{FC}$，称 $\mathbf{f}_N$ 为\textbf{内力}（施加它物体合力旋量为零）；若进一步 $\mathbf{f}_N\in\ker\mathbf{G}\cap\intop(\mathrm{FC})$（落在摩擦锥\emph{内部}），称\textbf{严格内力}。
\end{definition}

\begin{proposition}[力封闭 $\Leftrightarrow$ 满射 + 存在严格内力]
\label{prop:gc-internal-necessity}
抓取力封闭，当且仅当 $\mathbf{G}$ 满射（行满秩 $\rank\mathbf{G}=p$）\emph{且}存在严格内力 $\mathbf{f}_N\in\ker\mathbf{G}$ 使 $\mathbf{f}_N\in\intop(\mathrm{FC})$。
\end{proposition}
\begin{proof}
（充分）取任意 $\mathbf{w}_o\in\mathbb{R}^p$。$\mathbf{G}$ 满射故存在 $\mathbf{f}'$（未必在 $\mathrm{FC}$ 内）使 $\mathbf{w}_o=\mathbf{G}\mathbf{f}'$。考察 $\mathbf{f}'+\alpha\mathbf{f}_N$：当 $\alpha\to\infty$，$(\mathbf{f}'+\alpha\mathbf{f}_N)/\alpha\to\mathbf{f}_N\in\intop(\mathrm{FC})$，故存在足够大的 $\alpha'$ 使 $\mathbf{f}'+\alpha'\mathbf{f}_N\in\intop(\mathrm{FC})\subseteq\mathrm{FC}$，且 $\mathbf{G}(\mathbf{f}'+\alpha'\mathbf{f}_N)=\mathbf{G}\mathbf{f}'=\mathbf{w}_o$。故 $\mathbf{G}(\mathrm{FC})=\mathbb{R}^p$，力封闭。（必要）设力封闭。取 $\mathbf{f}_1\in\intop(\mathrm{FC})$ 使 $\mathbf{w}_o=\mathbf{G}\mathbf{f}_1\ne\mathbf{0}$；由力封闭取 $\mathbf{f}_2\in\mathrm{FC}$ 使 $\mathbf{G}\mathbf{f}_2=-\mathbf{w}_o$，令 $\mathbf{f}_N=\mathbf{f}_1+\mathbf{f}_2$，则 $\mathbf{G}\mathbf{f}_N=\mathbf{0}$ 且由锥性质 $\mathbf{f}_N\in\intop(\mathrm{FC})$，即严格内力。$\square$
\end{proof}

\noindent 直觉：力封闭抓取必须能“朝内夹”——存在一组净合力为零、却把各接触都压在摩擦锥内部的接触力。正是这股内力让各接触\emph{留有余量}：施加任意外力时，把外力对应的接触力叠加到内力上仍不越出摩擦锥，从而不打滑。无内力余量的抓取（接触力恰在锥面上）经不起任何切向外扰。

\subsection{构造性判据：正张成、凸包、分离超平面}
\label{sec:gc-fc-constructive}

直接验 $\mathbf{G}(\mathrm{FC})=\mathbb{R}^p$ 因摩擦锥形状而困难。对\emph{无摩擦点接触}（接触力只沿法向、$\mathbf{f}_{c_i}=f_i\hat{\mathbf{n}}_i$，$f_i\ge0$），$\mathbf{G}(\mathrm{FC})=\mathbb{R}^p$ 化为“$\mathbf{G}$ 各列的\emph{非负}线性组合张满 $\mathbb{R}^p$”，即正张成问题，由此得三条等价的可计算判据。先备定义\cite{murray1994mathematical}：向量组 $\{\mathbf{v}_i\}$ \textbf{正张成} $\mathbb{R}^p$，若任意 $\mathbf{x}\in\mathbb{R}^p$ 都可写成 $\sum_i a_i\mathbf{v}_i$（$a_i\ge0$）；集合 $S$ 的\textbf{凸包} $\conv(S)$ 是含 $S$ 的最小凸集；过点 $\mathbf{c}$、法向 $\mathbf{v}$ 的超平面把空间分两侧，若凸集 $K$ 全落一侧（$\mathbf{v}^\top(\mathbf{x}-\mathbf{c})\ge0,\forall\mathbf{x}\in K$）则称\textbf{分离超平面}。

\begin{theorem}[无摩擦点接触力封闭的凸性等价判据]
\label{thm:gc-fc-convex}
设无摩擦点接触抓取的抓取矩阵为 $\mathbf{G}\in\mathbb{R}^{p\times N}$、其列为 $\{\mathbf{G}_i\}$。下列陈述等价\cite{murray1994mathematical}：
\begin{enumerate}
  \item 抓取力封闭；
  \item $\mathbf{G}$ 的列 $\{\mathbf{G}_i\}$ 正张成 $\mathbb{R}^p$；
  \item 凸包 $\conv(\{\mathbf{G}_i\})$ 含原点的一个邻域（即原点在凸包\emph{内部}）；
  \item \emph{不}存在非零向量 $\mathbf{v}\in\mathbb{R}^p$ 使 $\mathbf{v}^\top\mathbf{G}_i\ge0$ 对一切 $i$ 成立（即不存在过原点的分离超平面）。
\end{enumerate}
\end{theorem}

\begin{note}[判据四为何最便于计算]
判据四把“是否力封闭”变成“能否找到一个把所有 $\mathbf{G}_i$ 压在同一侧的超平面法向 $\mathbf{v}$”：找得到 $\Rightarrow$ 非力封闭（$\mathbf{v}$ 方向的外力没人能抵抗），找不到 $\Rightarrow$ 力封闭。候选 $\mathbf{v}$ 不必盲搜——取 $\mathbf{G}$ 任意 $p-1$ 个独立列，其公共法向（正交补）就是一个候选 $\mathbf{v}$；只需逐个候选检查“$\mathbf{v}$ 与其余各列点积是否同号”。这是个有限枚举，$O\binom{N}{p-1}$ 量级，是无摩擦平面/空间抓取力封闭的标准手算/编程判据。$\mathbf{v}$（若存在）即过原点的支撑超平面法向。
\end{note}

\begin{example}[平面四点接触：聚拢于角 vs 共线退化]
\label{ex:gc-planar-fc}
矩形物体、四个无摩擦点接触\cite{murray1994mathematical}。\textbf{情形 A}（接触聚于四角，参数 $a,b>0$）：抓取矩阵
\[
  \mathbf{G}=\begin{bmatrix}1&0&-1&0\\ 0&-1&0&1\\ -a&b&-a&b\end{bmatrix}
  \quad(\text{平面 wrench }[f_x,f_y,m_z]^\top).
\]
枚举所有“两列定一支撑超平面”的候选 $\mathbf{v}_{ij}$，逐一验证均\emph{不}满足判据四的“同号”条件，故\textbf{力封闭}；其凸包含原点邻域。\textbf{情形 B}（接触取 $a,b=0$，法向都过几何中心）：
\[
  \mathbf{G}=\begin{bmatrix}1&0&-1&0\\ 0&-1&0&1\\ 0&0&0&0\end{bmatrix},
\]
第三行（力矩行）全零——无法用正组合抵抗任何力矩（判据二失效）；凸包整体落在 $f_x\text{-}f_y$ 平面、不含原点邻域（判据三失效）；取 $\mathbf{v}=[0,0,1]^\top$ 得 $\mathbf{v}^\top\mathbf{G}=[0,0,0,0]$，存在分离超平面（判据四命中）。故\textbf{非力封闭}：物体可绕中心自由转动。此例点睛——\emph{接触法向必须“张开力臂”，否则封不住转动自由度}。
\end{example}

\subsection{摩擦点接触：把摩擦锥的棱当无摩擦接触}
\label{sec:gc-fc-friction}

带摩擦时，摩擦锥内任一接触力都是锥\emph{棱}（edge）的非负组合。于是只需检查“与锥棱对应的物体 wrench 是否正张成 wrench 空间”——这把摩擦点接触归约成若干无摩擦点接触，从而可复用\cref{thm:gc-fc-convex}。平面 PCWF 接触的（每接触两条锥棱）接触映射形如
\begin{equation}
  \tilde{\mathbf{G}}_i=\begin{bmatrix}\hat{\mathbf{e}}_i^{+} & \hat{\mathbf{e}}_i^{-}\\ \mathbf{r}_i\times\hat{\mathbf{e}}_i^{+} & \mathbf{r}_i\times\hat{\mathbf{e}}_i^{-}\end{bmatrix},
  \qquad \mathrm{FC}_i'=\{\boldsymbol{\beta}\in\mathbb{R}^2:\beta_1,\beta_2\ge0\},
  \label{eq:gc-fc-edges}
\end{equation}
其中 $\hat{\mathbf{e}}_i^{\pm}$ 是第 $i$ 个摩擦锥两条棱的单位方向（半顶角 $\alpha=\arctan\mu$，\cref{def:friction-cone}）——形式与“两个独立的无摩擦点接触”完全相同。空间摩擦锥用多面（金字塔）线性化后同理（\cref{sec:contact-linearize} 的多棱锥）。

\begin{theorem}[Nguyen：空间三摩擦点接触力封闭的平面归约]
\label{thm:gc-nguyen}
设空间刚体被三个有摩擦点接触约束、三接触点\emph{不}共线，它们张成唯一平面 $S$。该抓取力封闭，当且仅当\emph{每个}接触的摩擦锥与 $S$ 相交于一个（二维）平面锥，\emph{且}平面 $S$ 上的这三个平面锥构成\textbf{平面力封闭}\cite{nguyen1988,lynch2017modern}。
\end{theorem}
\begin{proof}[证明思路]
（必要）若空间力封闭则其切片 $S$ 必平面力封闭；且若某摩擦锥与 $S$ 仅交于线或点，则绕“另两接触连线”的外力矩无法抵抗。（充分）取参考系使 $S$ 落于 $\hat{\mathbf{x}}\text{-}\hat{\mathbf{y}}$ 平面、$\hat{\mathbf{z}}$ 为法向 $N$。把各接触力与外力旋量正交分解为 $S$ 内分量与 $N$ 分量：$N$ 方向（法向力与面内力矩）由“三点不共线”保证三元线性方程组唯一可解；$S$ 方向（面内力与法向力矩）由 $S$ 的平面力封闭保证可解，且面内还可叠加平面内力使各接触力落入各自摩擦锥。两方向合成即得空间任意外扰的可行接触力。$\square$
\end{proof}

\begin{insight}[摩擦“买”到了什么：从 $\ge7$ 降到 $3$]
对照\cref{sec:gc-formclosure} 的形封闭——空间无摩擦点接触锁死物体要 $\ge7$ 个；而\cref{thm:gc-nguyen} 说空间\emph{有摩擦}只需 $3$ 个点接触即可力封闭。差距巨大的原因：无摩擦每接触只贡献 $1$ 维（法向）正向力，要正张成 $6$ 维 wrench 空间至少 $7$ 条射线（“正张成 $\mathbb{R}^n$ 需 $\ge n+1$ 个向量”）；有摩擦每接触贡献一个 $3$ 维摩擦锥（一束方向），$3$ 个锥（$3\times2=6$ 乃至更多锥棱在线性化后）足以正张成。\emph{摩擦把“一根针”换成了“一把扇面”，单接触的施力能力从一维升到三维（锥），于是接触数需求骤降}。代价是这份能力依赖 $\mu$——摩擦不足（如\cref{thm:gc-nguyen} 锥与 $S$ 退化成线）则失效，且接触力须始终保持在锥内（靠内力余量，\cref{prop:gc-internal-necessity}）。
\end{insight}

\begin{note}[力封闭与 GIWC 静平衡判据是同一族凸锥包含问题]
\cref{thm:giwc-lp} 的站立平衡判据“$-\mathbf{w}^{\text{gravity}}\in\mathcal{W}_{\text{GIWC}}$”与本节力封闭判据“$\mathbf{G}(\mathrm{FC})=\mathbb{R}^p$”同源：都把“各接触摩擦锥（$V$-表示，棱/生成元）经几何变换 Minkowski 求和张成的锥”作为可行 wrench 集，再问一个\emph{凸锥包含}问题。差别仅在\emph{问什么}：GIWC 问“某\emph{特定} wrench（重力反力）是否在锥内”（站不站得稳，一次 LP 可行性，\cref{eq:giwc}）；力封闭问“锥是否\emph{等于整个空间}”（能否抵抗\emph{任意}外扰，等价于原点在锥内部 / 锥的对偶退化为零）。$H/V$-表示对偶（\cref{thm:minkowski-weyl}）、Minkowski 和、双重描述法这套机器（\cref{sec:contact-cone-giwc}）对两者通用。\emph{力封闭可视为 GIWC 判据对“全方向外扰”的极端要求版}。
\end{note}

% ════════════════════════════════════════════════════════════════
\section{形封闭 Form Closure}
\label{sec:gc-formclosure}

\paragraph{动机：不靠摩擦也能锁死。}
形封闭是更强、更“放心”的封闭：一组\emph{静止}约束（手指/夹具/工装）从纯几何上封死物体的\emph{一切}运动，\emph{完全不依赖摩擦}。装配工装定位、零摩擦假设下的稳健夹持都要它——因为摩擦系数会变（油污、磨损），而几何锁死不会。形封闭正是\cref{sec:gc-multicontact} 那句“可行运动旋量锥缩成原点”的命名。

\subsection{定义与接触数定理}
\label{sec:gc-fc-form-def}

\begin{definition}[（一阶）形封闭]
\label{def:gc-form-closure}
若一组\emph{静止}约束使物体的可行运动旋量集 \cref{eq:gc-twist-cone} 仅含零旋量，则物体处于\textbf{形封闭}；约束由手指提供时称\textbf{形封闭抓取}。一阶形封闭指该结论由一阶分析（\cref{def:gc-first-order}，只用接触法向）得出。
\end{definition}

每个静止接触 $i$ 贡献半空间约束 $\mathcal{F}_i^\top\mathcal{V}\ge0$。形封闭要求满足全部约束的唯一旋量是 $\mathcal{V}=\mathbf{0}$，对 $j$ 个接触等价于\emph{接触 wrench 正张成全空间}：
\begin{equation}
  \pos(\{\mathcal{F}_1,\dots,\mathcal{F}_j\})=\mathbb{R}^p
  \quad(\text{空间 }p=6\text{；平面 }p=3).
  \label{eq:gc-form-posspan}
\end{equation}

\begin{theorem}[一阶形封闭的接触数下界]
\label{thm:gc-contact-count}
平面物体一阶形封闭至少需 $4$ 个点接触；空间物体至少需 $7$ 个点接触\cite{lynch2017modern}.
\end{theorem}
\begin{proof}
正张成 $\mathbb{R}^n$ 的向量组至少含 $n+1$ 个向量（“$n$ 个正交基加上其负和” $-\sum_i\mathbf{a}_i$ 给出恰 $n+1$ 个正张成 $\mathbb{R}^n$ 的向量；少于 $n+1$ 个则必存在某方向 $\mathbf{c}$ 使所有 $\mathbf{a}_i^\top\mathbf{c}\le0$，留出未封死的运动）。无摩擦点接触每个贡献一个法向力旋量，故平面 $p=3$ 需 $3+1=4$、空间 $p=6$ 需 $6+1=7$。$\square$
\end{proof}

\begin{note}[例外物体：圆盘永远封不住]
\cref{thm:gc-contact-count} 有个反例类——\textbf{例外物体}（exceptional objects）：其所有接触点法向力的正张成永远 $\ne\mathbb{R}^n$，无论接触多少个都无法形封闭。平面圆盘是典型——任意位置法向都过圆心，绕圆心的转动永远封不住；空间里旋转面（球、椭球）同理。这类物体只能靠\emph{摩擦}（力封闭）或\emph{二阶几何}（曲率）锁住，纯一阶形封闭对它们束手无策。
\end{note}

\begin{note}[一阶充分、非必要：二阶曲率可减接触数]
\cref{thm:gc-contact-count} 是\emph{一阶}下界。一阶分析只看法向、不看曲率，故保守：某些抓取一阶判为“可绕某中心转动”（非形封闭），但二阶分析（计入接触曲率 $\partial^2 d/\partial\mathbf{q}^2$）显示该转动实际穿透或分离、物体其实被锁死。经典例：平面两指若指尖与物体曲率配合得当，\emph{两}个接触即可二阶形封闭，远少于一阶要求的四个\cite{lynch2017modern}。一阶判据的安全方向是单边的——\emph{一阶判为形封闭则任意阶都形封闭}（铁结论）；\emph{一阶判为仅 roll-slide 可行则可能被高阶分析救成形封闭}；一阶判为有分离/穿透运动则任意阶都非形封闭。
\end{note}

\subsection{线性规划判据}
\label{sec:gc-fc-form-lp}

\begin{theorem}[一阶形封闭的 LP 判据]
\label{thm:gc-form-lp}
设 $\mathbf{F}=[\mathcal{F}_1\ \mathcal{F}_2\ \cdots\ \mathcal{F}_j]\in\mathbb{R}^{p\times j}$ 列为各接触 wrench（空间 $p=6$；平面 $p=3$，$\mathcal{F}_i=[f_{ix},f_{iy},m_{iz}]^\top$，本书序力在前）。物体一阶形封闭 $\Leftrightarrow$ $\rank\mathbf{F}=p$ \emph{且}下列线性规划有解\cite{lynch2017modern}：
\begin{equation}
  \begin{aligned}
    &\text{求}\ \mathbf{k}\in\mathbb{R}^j,\\
    &\min\ \mathbf{1}^\top\mathbf{k}\\
    &\text{s.t.}\ \ \mathbf{F}\mathbf{k}=\mathbf{0},\quad k_i\ge1,\ i=1,\dots,j,
  \end{aligned}
  \label{eq:gc-form-lp}
\end{equation}
$\mathbf{1}$ 为全 $1$ 向量。若 $\mathbf{F}$ 满秩且 \cref{eq:gc-form-lp} 有解，则一阶形封闭；否则不是。
\end{theorem}

\noindent 解读：$\rank\mathbf{F}=p$ 保证接触 wrench 张成（线性意义）全空间；约束 $\mathbf{F}\mathbf{k}=\mathbf{0},\ k_i\ge1$（严格正系数）保证存在“各接触都\emph{真正}出力、合力旋量为零”的内力配置——这正是 \cref{eq:gc-form-posspan} 正张成的代数判定（原点在 wrench 凸包\emph{内部}而非边界）。目标 $\mathbf{1}^\top\mathbf{k}$ 对二值判定非必需，仅令 LP 良定。

\begin{example}[内孔双指抓取的 LP 判定]
\label{ex:gc-form-lp}
平面物体中央有孔，两指各触孔的两条边，产生四个接触法向\cite{lynch2017modern}。本书序（力在前 $[f_x,f_y,m_z]$）的
\[
  \mathbf{F}=\begin{bmatrix}-1&0&1&0\\ 0&1&0&-1\\ 0&0&2&-2\end{bmatrix}\ (\rank=3),
\]
$\mathbf{F}$ 的零空间一维、由 $[1,1,1,1]^\top$ 张成，LP \cref{eq:gc-form-lp} 返回 $k_1=k_2=k_3=k_4=1$（满足 $\mathbf{F}\mathbf{k}=\mathbf{0}$、$k_i\ge1$），故\textbf{形封闭}。若把右指移到孔的右下角，$\mathbf{F}$ 仍满秩但 LP 无解——\emph{非}形封闭，物体可沿页面向下滑出。这给出工装/夹具设计的可计算闸门：调接触位置直到 LP 可行。
\end{example}

\subsection{抓取质量度量}
\label{sec:gc-fc-quality}

两个都形封闭的抓取，谁更好？需\textbf{抓取质量度量} $\mathrm{Qual}(\{\mathcal{F}_i\})$：$<0$ 表非形封闭，正值越大越好。一种常用度量：限定每接触力幅 $f_i\le f_{i,\max}$，则各接触可施加的总 wrench 集为凸多胞形
\begin{equation}
  C_F=\Big\{\sum_{i=1}^{j}f_i\mathcal{F}_i:\ f_i\in[0,f_{i,\max}]\Big\},
  \label{eq:gc-quality-polytope}
\end{equation}
形封闭时 $C_F$ 含原点于内部。取 $\mathrm{Qual}=$ “以 wrench 空间原点为心、可内接于 $C_F$ 的最大球半径”（即原点到 $C_F$ 各边界超平面的最小有符号距离）\cite{ferrari1992,lynch2017modern}。两处须留意：（1）力与力矩量纲不同，须用物体特征长度 $r$ 把力矩 $m$ 折成 $m/r$ 再比；（2）力矩依赖参考原点选取，通常取物体几何中心或质心。这个“最大内接球半径”度量（即著名的 Ferrari–Canny 度量）也直接用于力封闭抓取的打分与优选。

% ════════════════════════════════════════════════════════════════
\section{形封闭 vs 力封闭：对比与抉择}
\label{sec:gc-compare}

\begin{table}[H]\centering\small
\caption{形封闭与力封闭对照}
\label{tab:gc-form-vs-force}
\begin{tabular}{p{2.6cm}p{5.4cm}p{5.4cm}}
\toprule
维度 & 形封闭 form closure & 力封闭 force closure \\
\midrule
核心定义 & 静止约束几何锁死物体一切运动：可行旋量锥 $=\{\mathbf{0}\}$（\cref{def:gc-form-closure}） & 接触力能抵抗任意外力旋量：$\mathbf{G}(\mathrm{FC})=\mathbb{R}^p$（\cref{def:gc-force-closure}） \\
依赖摩擦？ & \textbf{否}（纯几何，$\mu=0$ 也成立） & \textbf{是}（靠摩擦锥；$\mu=0$ 时退化为一阶形封闭） \\
等价判据 & wrench 正张成 $\mathbb{R}^p$ / 原点在 wrench 凸包内部 / LP 可行（\cref{thm:gc-form-lp}） & wrench 正张成 / 原点在凸包内部 / 无分离超平面（\cref{thm:gc-fc-convex}） \\
接触数（空间） & $\ge7$ 点接触（\cref{thm:gc-contact-count}） & $\ge3$ 摩擦点接触（不共线，\cref{thm:gc-nguyen}） \\
接触数（平面） & $\ge4$ 点接触 & $\ge2$ 摩擦点接触 \\
判定工具 & LP（\cref{eq:gc-form-lp}） & LP / 凸包 / 分离超平面 / Nguyen 平面归约 \\
强弱关系 & \textbf{更强}：形封闭 $\Rightarrow$ 力封闭 & \textbf{更弱}：力封闭 $\not\Rightarrow$ 形封闭 \\
典型用途 & 工装/夹具定位、零摩擦稳健夹持 & 灵巧手抓取、少指抓持、装配夹持 \\
盲区/代价 & 例外物体（圆盘）封不住；接触数多、笨重；一阶保守（曲率可救） & 依赖 $\mu$（油污/磨损失效）；须内力余量防滑（\cref{prop:gc-internal-necessity}） \\
与本书锚点 & \cref{sec:gc-multicontact} 旋量锥；\cref{thm:giwc-lp} 同族 LP & \cref{ch:mr-coopforce} 抓取映射 $\mathbf{G}$；\cref{def:friction-cone} 摩擦锥；\cref{sec:contact-cone-giwc} GIWC \\
\bottomrule
\end{tabular}
\end{table}

\begin{insight}[一张图记住二者关系：形封闭 $\subsetneq$ 力封闭]
最该刻进脑子的一句：\emph{形封闭是力封闭在 $\mu=0$（且改用法向单向力正张成）下的特例的反面强化——形封闭蕴含力封闭，反之不然}。几何上，形封闭把可行\emph{运动}锥压成原点（运动侧封死），力封闭把可施\emph{力}锥撑满全空间（力侧满射）——由运动-静力学对偶（\cref{eq:cm-duality}），“运动锥 $=\{\mathbf{0}\}$”与“力锥 $=\mathbb{R}^p$”本是同一枚硬币两面，只是形封闭用\emph{无摩擦法向力}撑、力封闭用\emph{摩擦锥}撑，后者“扇面”更宽故更易满射、所需接触更少。工程取舍：要\emph{稳健}（不信摩擦）选形封闭、认接触多；要\emph{灵巧/省指}选力封闭、认须管摩擦与内力。
\end{insight}

% ════════════════════════════════════════════════════════════════
\section*{常见误解汇总}
\begin{enumerate}
  \item \textbf{“力封闭就等于物体不会动。”}——错。力封闭只保证摩擦锥能\emph{生成}任意 wrench（能力），不保证电机\emph{已经}把物体顶住；是否静止还取决于内力大小与电机出力方向（\cref{def:gc-force-closure} 后的 note）。
  \item \textbf{“滚动接触一定有转动。”}——错。滚动只指接触点相对速度为零（\cref{eq:gc-rolling}），含两体相对静止（粘连）的情形；区分 R/S 看切向有无滑移，不看整体转动。
  \item \textbf{“接触越多越容易形封闭。”}——半对。接触数是\emph{必要}下界（平面 $\ge4$、空间 $\ge7$），但例外物体（圆盘）无论多少接触都封不住；且接触位置退化（共线/法向全过中心，\cref{ex:gc-planar-fc} 情形 B）会让多接触仍失效。
  \item \textbf{“$\mathbf{J}_h$ 和 $\mathbf{G}^\top$ 是一回事。”}——错。$\mathbf{G}^\top$ 是物体运动到接触点速度的纯几何分发（只含 $\mathbf{r}_i$），$\mathbf{J}_h$ 是手指关节到指尖速度的手指运动学（含 $\boldsymbol{\theta}$）；基本抓取约束要二者在可施力方向相等（\cref{thm:gc-grasp-constraint} 后的 note）。
  \item \textbf{“力封闭判据可以不管手指自由度。”}——半对。$\mathbf{G}(\mathrm{FC})=\mathbb{R}^p$ 只看接触几何与摩擦，确实不含 $\mathbf{J}_h$；但其隐含前提是“$\mathrm{FC}$ 内的力都\emph{施得出}”，须 $\mathbf{J}_h$ 行满秩、无不可控结构力（\cref{prop:gc-structural}）。手不够灵巧时，力封闭的接触配置也未必可控。
  \item \textbf{“一阶形封闭判否就一定封不住。”}——错。一阶保守：一阶判“仅 roll-slide 可行”时，二阶曲率分析可能把物体救成形封闭（如曲率配合的两指，\cref{sec:gc-fc-form-def} 的 note）；但一阶判“形封闭”则任意阶都形封闭（这一方向是铁的）。
  \item \textbf{“wrench 排序无所谓。”}——半对。结论（功、对偶、张成性、封闭性）确与排序无关；但\emph{抄公式}时必须对齐排序——本书力在前 $[\mathbf{f};\mathbf{m}]$，MR/MLS 力矩在前 $(\mathbf{m},\mathbf{f})$，混用会让 $\mathbf{F}$ 矩阵行块错位、LP 解错（\cref{eq:gc-normal-wrench} 的 note）。
\end{enumerate}

\section*{本章小结}
本章把\textbf{抓取映射} $\mathbf{G}$（\cref{ch:mr-coopforce}）与\textbf{接触摩擦锥}（\cref{ch:contact-mechanics}）合流，建立了抓取封闭性的完整理论链：
\begin{enumerate}
  \item \textbf{接触一阶运动学}（\cref{sec:gc-contactkin}）：距离函数 $\dot d=0$ 给不可穿透/激活约束 $\mathcal{F}^\top(\mathcal{V}_A-\mathcal{V}_B)\gtrless0$；据接触点相对速度分滚动 R / 滑动 S / 分离 B；多接触可行旋量集是齐次多面体凸锥 \cref{eq:gc-twist-cone}。
  \item \textbf{接触力旋量基}（\cref{sec:gc-wrenchbasis}）：无摩擦点接触 $\mathbf{B}_{c_i}$ $1$ 维、PCWF $3$ 维、软指 SF $4$ 维；摩擦锥 $\mathrm{FC}$ 是闭凸锥；$\mathbf{G}_i=\mathrm{Ad}_{g_{oc_i}^{-1}}^\top\mathbf{B}_{c_i}$ 接回\cref{eq:cm-Gi}。
  \item \textbf{手部雅可比}（\cref{sec:gc-handjac}）：基本抓取约束 $\mathbf{J}_h\dot{\boldsymbol{\theta}}=\mathbf{G}^\top\mathcal{V}$；可抓取 $\Leftrightarrow\range(\mathbf{G}^\top)\subseteq\range(\mathbf{J}_h)$；结构力 $\in\ker\mathbf{J}_h^\top$。两 SCARA 手指抓盒的贯通数值范例（\cref{sec:gc-scara}）实证了\emph{力封闭却不可抓取}、并解出 2 维内部运动。
  \item \textbf{力封闭}（\cref{sec:gc-forceclosure}）：$\mathbf{G}(\mathrm{FC})=\mathbb{R}^p$；三等价判据（正张成 / 原点在凸包内部 / 无分离超平面）；须严格内力（\cref{prop:gc-internal-necessity}）；Nguyen 三接触平面归约（\cref{thm:gc-nguyen}）；与 GIWC 同族。
  \item \textbf{形封闭}（\cref{sec:gc-formclosure}）：纯几何、可行旋量锥 $=\{\mathbf{0}\}$；平面 $\ge4$ / 空间 $\ge7$ 接触（\cref{thm:gc-contact-count}）；LP 判据 \cref{eq:gc-form-lp}；一阶充分非必要、例外物体封不住。
  \item \textbf{对比}（\cref{sec:gc-compare}）：形封闭 $\Rightarrow$ 力封闭、反之不然；稳健选形封闭、灵巧选力封闭。
\end{enumerate}
封闭性判据本身只回答“给定接触能否抓稳”；如何\emph{选}接触点使之封闭、并据质量度量打分优选（乃至数据驱动学习抓取位姿），属\emph{抓取规划}范畴，与本章互为“判据—综合”的两面，建议接续研读相关抓取规划与抓取合成文献（见\nameref{ch:notation} 后的延伸阅读）。

\section*{速查卡}
\begin{table}[H]\centering\small
\caption{核心概念速查}
\label{tab:gc-cheatsheet}
\begin{tabular}{cll}
\toprule
\# & 概念 & 一句话 / 关键式 \\
\midrule
1 & 不可穿透约束 & $\mathcal{F}^\top(\mathcal{V}_A-\mathcal{V}_B)\ge0$；$=0$ 接触维持、$>0$ 分离 \\
2 & R/S/B 三模式 & 接触点相对速度：$=0$ 滚动；切向滑移滑动；分离 breaking \\
3 & 力旋量基 $\mathbf{B}_{c_i}$ & 无摩擦 $1$ 维 / PCWF $3$ 维 / 软指 $4$ 维 \\
4 & 摩擦锥 & 闭凸锥；$\sqrt{f_x^2+f_y^2}\le\mu f_z$，半顶角 $\arctan\mu$ \\
5 & 抓取映射 $\mathbf{G}$ & $\mathbf{w}=\mathbf{G}\mathbf{f}$，$\mathbf{G}_i=[\mathbf{I};\skewm{\mathbf{r}}_i]$（PCWF） \\
6 & 基本抓取约束 & $\mathbf{J}_h\dot{\boldsymbol{\theta}}=\mathbf{G}^\top\mathcal{V}$ \\
7 & 可抓取性 & $\range(\mathbf{G}^\top)\subseteq\range(\mathbf{J}_h)$ \\
8 & 结构力 & $\in\ker\mathbf{J}_h^\top$，电机不可控 \\
9 & 力封闭 & $\mathbf{G}(\mathrm{FC})=\mathbb{R}^p$；$\Leftrightarrow$ 满射 $+$ 严格内力 \\
10 & 力封闭判据 & 正张成 / 原点在凸包内部 / 无分离超平面 \\
11 & Nguyen 定理 & 空间 $3$ 不共线摩擦点 $\Rightarrow$ 力封闭（锥交平面、$S$ 平面力封闭） \\
12 & 形封闭 & 可行旋量锥 $=\{\mathbf{0}\}$；纯几何无摩擦 \\
13 & 接触数下界 & 平面 $\ge4$、空间 $\ge7$ 点接触（一阶） \\
14 & 形封闭 LP & $\rank\mathbf{F}=p$ 且 $\mathbf{F}\mathbf{k}=\mathbf{0},k_i\ge1$ 有解 \\
15 & 抓取质量 & $C_F$ 内接最大球半径（Ferrari–Canny，力矩折 $m/r$） \\
16 & 强弱关系 & 形封闭 $\Rightarrow$ 力封闭；反之不然 \\
\bottomrule
\end{tabular}
\end{table}

\begin{table}[H]\centering\small
\caption{符号表}
\label{tab:gc-symbols}
\begin{tabular}{lll}
\toprule
符号 & 含义 & 首次出现 \\
\midrule
$d(\mathbf{q})$ & 两体距离函数（$>0$ 分离、$=0$ 相切、$<0$ 穿透） & \cref{eq:gc-d-derivs} \\
$\hat{\mathbf{n}}$ & 接触法向单位向量（指向物体 $A$ 内） & \cref{eq:gc-dot-pn} \\
$\mathcal{V}=[\mathbf{v};\boldsymbol{\omega}]$ & 运动旋量（twist，平移在前） & \cref{sec:gc-impenetrability} \\
$\mathcal{F}=[\mathbf{f};\mathbf{m}]$ & 力旋量（wrench，本书序力在前） & \cref{eq:gc-normal-wrench} \\
$\mathbf{B}_{c_i}\in\mathbb{R}^{p\times m_i}$ & 第 $i$ 接触力旋量基 & \cref{def:gc-contact-models} \\
$\mathrm{FC}_{c_i}$ & 第 $i$ 接触摩擦锥（闭凸锥） & \cref{eq:gc-Bc-pcwf} \\
$\mu,\gamma$ & 切向 / 扭转摩擦系数 & \cref{eq:gc-Bc-sf} \\
$\mathbf{G}\in\mathbb{R}^{p\times m}$ & 抓取映射，$\mathbf{w}=\mathbf{G}\mathbf{f}_c$ & \cref{eq:gc-graspmap} \\
$\mathbf{r}_i$ & 接触点相对物体质心位置 & \cref{eq:gc-graspmap} \\
$\mathbf{J}_h\in\mathbb{R}^{m\times n}$ & 手部雅可比 & \cref{eq:gc-Jh-def} \\
$\boldsymbol{\theta},\dot{\boldsymbol{\theta}}$ & 手指关节角 / 速度（总维 $n$） & \cref{eq:gc-grasp-constraint} \\
$\mathbf{f}_N$ & 内力（$\in\ker\mathbf{G}\cap\mathrm{FC}$） & \cref{def:gc-internal} \\
$\mathbf{F}=[\mathcal{F}_1\cdots\mathcal{F}_j]$ & 接触 wrench 矩阵（形封闭 LP） & \cref{thm:gc-form-lp} \\
$\pos,\conv,\intop$ & 正张成 / 凸包 / 内部 & \cref{thm:gc-fc-convex} \\
\bottomrule
\end{tabular}
\end{table}

\section*{故障排查手册}
\begin{enumerate}
  \item \textbf{症状}：力封闭判据时而真时而假、结果飘忽。\textbf{排查}：wrench 排序是否一致？本书力在前 $[\mathbf{f};\mathbf{m}]$，若混入 MR/MLS 的力矩在前 $(\mathbf{m},\mathbf{f})$，$\mathbf{F}$ 行块错位（\cref{eq:gc-normal-wrench} note）。统一排序后重算。
  \item \textbf{症状}：明明七个接触，形封闭 LP 仍无解。\textbf{原因}：物体是例外物体（圆盘/旋转面），或接触退化（共线、法向全过几何中心，\cref{ex:gc-planar-fc} B）。\textbf{对策}：核对 $\rank\mathbf{F}=p$；非例外物体则移动接触点打破退化。
  \item \textbf{症状}：力封闭判定通过，实物却打滑掉落。\textbf{原因}：力封闭是“能力”非“已实现平衡”（\cref{def:gc-force-closure} note）；内力余量不足或电机出力方向受限。\textbf{对策}：加大内力 $\mathbf{f}_{\text{int}}$（\cref{eq:cm-decomp}），核对各接触力是否留在锥内部（\cref{prop:gc-internal-necessity}）。
  \item \textbf{症状}：手部雅可比算出的可抓取性与直觉矛盾。\textbf{原因}：把 $\mathbf{G}^\top$ 当 $\mathbf{J}_h$、或漏了 $\mathbf{B}_{c_i}^\top$ 的屏蔽（\cref{eq:gc-Jh-def}）。\textbf{对策}：核对 $\mathbf{J}_h$ 含手指构型 $\boldsymbol{\theta}$、$\mathbf{G}$ 只含 $\mathbf{r}_i$；分别验 $\range(\mathbf{G}^\top)$ 与 $\range(\mathbf{J}_h)$ 的包含。
  \item \textbf{症状}：滚动/滑动判断与摩擦力方向对不上。\textbf{原因}：把“滚动 $=$ 有转动”当判据。\textbf{对策}：只看接触点相对速度（\cref{eq:gc-rolling}）——为零即滚动（含粘连），切向滑移即滑动；据此再定摩擦力（静摩擦取锥内、动摩擦反滑移）。
  \item \textbf{症状}：力封闭分析说三接触够、实测需更多。\textbf{原因}：Nguyen 定理要求摩擦锥与接触平面 $S$ 交于\emph{二维}锥（\cref{thm:gc-nguyen}）；$\mu$ 不足时锥退化成线/点，定理前提失效。\textbf{对策}：提高 $\mu$ 或增设接触；检查三点是否近共线（绕其连线的力矩无人抵抗）。
\end{enumerate}
